\documentclass[12pt]{article}
\usepackage[utf8]{inputenc}
\usepackage[T1]{fontenc}
\usepackage{hyperref}
\title{Uamu\'e Language Resources for K-12 Curriculum Development}
\author{Compiled by Claude Sonnet 5 (claude-sonnet-5) for Andr\'as Kornai}
\date{}
\begin{document}
\maketitle
\section{Introduction}
Uamu\'e (ISO 639-3 uam; also spelled Huamu\^e, Wamo\'e) is a dormant or extinct language of Brazil. Sources disagree on its classification: Wikipedia calls it unclassified, with speculation about a Kariri link, and Ethnologue calls it an isolate. A word list was reportedly recorded from a shaman of the Atikum near Cachoeirinha, Pernambuco. Do not confuse it with Guamo/Wamo, which is a different language. Very little material exists.
\section{Primary Educational Resources}
No primers, textbooks or school programmes were found in the searches.
\section{Linguistic and Reference Materials}
\begin{itemize}
\item Wikipedia, Uamu\'e language: \url{https://en.wikipedia.org/wiki/Uamu%C3%A9_language}
\item Wikipedia, List of unclassified languages of South America: \url{https://en.wikipedia.org/wiki/List_of_unclassified_languages_of_South_America}
\item Ethnologue entry for uam: \url{https://www.ethnologue.com/language/uam}
\end{itemize}
\section{Religious and Community Resources}
None found. The Atikum community is mentioned in the Wikipedia article, but no community organisation or contact was found.
\section{Digital and Audio Resources}
No audio or digital resources were found.
\section{Curriculum Development Context}
No native speakers remain and no teaching materials were found. Any curriculum would have to be a heritage-awareness programme built on the sparse word lists, ideally developed with the Atikum community. The sources above only point to the existence of those lists, and I did not examine them.
\end{document}